\chapter{Propulsión cohete: misiones y lenguaje físico}\label{ch:01}
\objetivos{Distinguir fuente de energía y masa de reacción; clasificar los sistemas propulsivos; relacionar los requisitos de una misión con empuje, impulso y potencia; separar dato, modelo y verificación.}
\section{Qué hace un sistema propulsivo}
Un motor cohete transfiere cantidad de movimiento a una masa que expulsa. El vehículo recibe la reacción correspondiente. La conservación se aplica al conjunto vehículo más materia expulsada; aplicada sólo a la parte que permanece a bordo requiere contabilizar el flujo de masa. Esta distinción explica por qué no es necesario empujar contra el aire y por qué un sistema puede funcionar en vacío. El entorno modifica presiones, intercambio térmico y fuerzas externas, pero no proporciona la reacción indispensable para la propulsión.

Conviene separar tres preguntas. La primera es qué sustancia abandona el vehículo: la masa de reacción. La segunda es de dónde procede la energía que permite acelerarla. La tercera es cómo se transforma esa energía en movimiento dirigido. En un motor químico las sustancias transportadas participan en reacciones que cambian su energía química. En un propulsor eléctrico, una fuente eléctrica aporta una parte esencial de la energía; el fluido puede calentarse, ionizarse o acelerarse mediante campos. La fuente eléctrica y el propulsante no son necesariamente el mismo subsistema.

La palabra \emph{propulsante} incluye la materia consumida por la propulsión. En un bipropulsante químico comprende combustible y oxidante; usar sólo el caudal de combustible en el denominador del impulso específico produciría una comparación falsa. Un monopropulsante no se define por almacenar combustible sin oxidante, sino por funcionar con una sustancia que puede liberar energía mediante un proceso adecuado sin un segundo propulsante separado. Un sólido puede contener componentes de distintas funciones distribuidos en una matriz; su clasificación no sustituye al análisis de su respuesta mecánica y térmica.

Las magnitudes de interés se definen antes de calcular. El empuje $F$ es una fuerza, en newtons. El impulso total es una integral temporal, en newton segundo. El caudal $\dot m$ es masa por tiempo. La velocidad efectiva $c=F/\dot m$ tiene dimensiones de velocidad aunque incluya un término de presión y no coincida con la velocidad física en la salida. El impulso específico $I_{sp}=c/g_0$ se expresa en segundos; $g_0=9.80665\,\mathrm{m\,s^{-2}}$ es una constante convencional, no la aceleración gravitatoria en la posición del vehículo.
\section{Clasificar para formular modelos}
Los motores químicos térmicos convierten energía interna y química en entalpía de los productos, y una tobera transforma parte de ella en energía cinética dirigida. La clasificación en sólidos, líquidos e híbridos informa sobre almacenamiento y mecanismos de suministro. No determina por sí sola el impulso específico, la estabilidad, el número de encendidos o el grado de regulación. Esas propiedades requieren examinar la arquitectura concreta y las condiciones de operación.

En un sólido, la superficie reactiva evoluciona al consumirse material. Su caudal no es una variable arbitraria: resulta del acoplamiento entre regresión, presión y descarga. En un líquido, las líneas de alimentación, la inyección, la combustión y la tobera forman una red acoplada. Las bombas aumentan presión a costa de potencia; el sistema que la proporciona participa en el balance completo. En un híbrido las fases de los propulsantes son distintas y la regresión puede estar asociada al flujo que recorre la superficie. No debe trasladarse automáticamente una ley empírica de un sólido a un híbrido.

La propulsión eléctrica comprende sistemas electrotérmicos, electrostáticos y electromagnéticos. Un Hall contiene campos magnéticos importantes para el transporte electrónico, pero la aceleración principal de los iones se asocia al campo eléctrico. Clasificarlo por la mera presencia de un imán ocultaría la física que se pretende estudiar. En todos los casos hay que cerrar un presupuesto de potencia, masa, impulso y tiempo de operación.

Las comparaciones requieren fronteras homogéneas. Una comparación de masas puede incluir sólo el propulsante, el sistema de alimentación, el motor seco o el subsistema completo con fuente de energía. Una comparación de rendimiento puede referirse a la energía del chorro, la conversión eléctrica o la utilidad para una maniobra. Ninguna de estas razones es una eficiencia universal del vehículo.
\section{Empuje, duración y utilidad de una misión}
La ecuación ideal del cohete relaciona una velocidad efectiva con una razón de masas. No dice cuánto tarda la maniobra. Si dos sistemas producen el mismo impulso total, no necesariamente llevan a cabo la misma trayectoria: durante tiempos distintos actúan fuerzas externas y cambian las condiciones orbitales. Tampoco el impulso total basta para comparar dos motores que aceleran vehículos con masas diferentes.

Para una descarga unidireccional ideal, sin presión adicional, se tiene
\begin{align}
 F&=\dot m v_e, & \dot E_j&=\frac12\dot m v_e^2,\\
 \frac{\dot E_j}{F}&=\frac{v_e}{2}, & F&=\frac{2\dot E_j}{v_e}.
\end{align}
La primera relación cuenta momento; la segunda cuenta energía cinética del chorro respecto al motor. A potencia de chorro fijada, aumentar la velocidad de escape reduce el empuje. No es una contradicción: se consume menos masa, pero cada kilogramo exige más energía. Esta relación será central en propulsión eléctrica. No se aplica sin modificación si se confunden potencia eléctrica de entrada y potencia dirigida del chorro.

En una misión espacial aparecen restricciones adicionales: tiempo disponible, orientación, masa de los depósitos, almacenamiento, potencia solar, disipación térmica y fiabilidad. Un sistema de alto $I_{sp}$ puede ahorrar propulsante y necesitar una fuente de potencia cuya masa elimine esa ventaja. Un motor de gran empuje puede reducir pérdidas gravitatorias y, a la vez, exigir estructura o alimentación más pesada. La decisión requiere un problema de sistema, no ordenar una tabla por un único número.
\section{De la historia a una forma de razonar}
Como hitos documentados, NASA sitúa el primer vuelo de un cohete de propulsante líquido de Robert Goddard el 16 de marzo de 1926, en Auburn, Massachusetts. El lanzamiento de Sputnik el 4 de octubre de 1957 abrió la etapa de satélites artificiales en órbita. Son hitos de demostración y de misión distintos: hacer funcionar una arquitectura propulsiva y alcanzar una condición orbital no exigen el mismo conjunto de prestaciones.\cite{NASA-GODDARD,NASA-SPUTNIK}

La evolución posterior puede entenderse mediante compromisos de sistema. Un lanzador necesita mucho empuje durante una fase breve; una sonda con energía eléctrica puede acumular un gran cambio de velocidad a lo largo de intervalos prolongados. En la historia de la propulsión eléctrica descrita por Goebel y Katz, el desarrollo de fuentes de plasma, aceleración y sistemas de potencia permitió pasar de demostraciones a aplicaciones de mantenimiento orbital y misiones científicas. No es una sustitución universal de la propulsión química: amplía el conjunto de soluciones disponibles para objetivos distintos.

La historia de la propulsión no consiste sólo en una sucesión de récords. El reconocimiento de la masa variable permitió separar la necesidad de transportar masa de reacción de la necesidad de una atmósfera. El desarrollo de toberas supersónicas hizo posible convertir una mayor fracción de entalpía en velocidad dirigida. Las arquitecturas líquidas añadieron capacidad de suministro y control; los sólidos aportaron otras formas de almacenamiento e integración; la propulsión eléctrica amplió el compromiso entre consumo y duración.

Estos avances no invalidaron las conservaciones. Cambiaron los mecanismos con los que se satisfacen y los límites de los modelos simplificados. Un mismo balance de momento permite estudiar una tobera química o un haz de iones, pero las ecuaciones constitutivas que suministran el caudal, la energía y las fuerzas internas son diferentes. Esa continuidad conceptual es más útil que memorizar una fórmula distinta para cada tecnología.

En este manual los seis bloques proceden de la documentación docente ULE; los 28 capítulos son una división editorial. La guía establece 6 ECTS y 150 horas de trabajo, con un 75\,\% de prueba escrita y un 25\,\% de trabajo y defensa. Los tiempos y puntuaciones de los simulacros son propuestas de estudio. La oferta de TFG y TFM acompaña a los materiales, pero no se utiliza como programa de la asignatura.\cite{GUIA2026,ULE-PRES}
\section{Problema resuelto: dos sistemas con distinta velocidad efectiva}
\begin{studiumworked}{M01. Cantidad de movimiento y presupuesto de potencia}
\textbf{Enunciado y datos.} Dos modelos ideales entregan la misma potencia al chorro, $\dot E_j=600\,\mathrm W$. Sus velocidades son $v_{e,A}=3000\,\mathrm{m/s}$ y $v_{e,B}=18000\,\mathrm{m/s}$. Determine empuje, caudal e impulso en un intervalo de $120\,\mathrm s$.

\textbf{Incógnitas e hipótesis.} $F$, $\dot m$ e $I$. Régimen constante, chorro axial monocinético, entorno sin término de presión y potencia ya referida al chorro. No se prescribe una tecnología física ni una eficiencia eléctrica.

\textbf{Modelo y despeje.} Conservar momento y energía da
\[F=2\dot E_j/v_e,\qquad \dot m=2\dot E_j/v_e^2,\qquad I=F\Delta t.\]
\textbf{Sustitución y resultados.} Para A: $F_A=1200/3000=0.400\,\mathrm N$, $\dot m_A=1.33333\times10^{-4}\,\mathrm{kg/s}$ e $I_A=48.0\,\mathrm{N\,s}$. Para B: $F_B=0.0666667\,\mathrm N$, $\dot m_B=3.70370\times10^{-6}\,\mathrm{kg/s}$ e $I_B=8.00\,\mathrm{N\,s}$.

\textbf{Interpretación.} Multiplicar por seis la velocidad divide el empuje por seis y el caudal por 36 a potencia fija. La masa ahorrada no se convierte automáticamente en una maniobra más rápida.

\textbf{Comprobación.} En ambos casos $Fv_e/2=600\,\mathrm W$ y $I=\Delta m\,v_e$. Las dimensiones de $2\dot E_j/v_e$ son newtons. \textbf{Límites.} Se excluyen divergencia, distribuciones de velocidad, pérdidas, masa de la fuente y variación de masa del vehículo.
\textbf{Operaciones para comparar ambos casos.}
\begin{enumerate}

\item Despejar empuje de potencia cinética: F=2P/v.

\[\fitmath{\mathrm{FA}=\frac{2\cdot 600}{3000}\approx 0.4\,\mathrm{N}}\]

\item Del balance de energía, caudal es 2P/v al cuadrado.

\[\fitmath{\mathrm{mA}=\frac{2\cdot 600}{{3000}^{2}}\approx 0.0001333333\,\mathrm{kg/s}}\]

\item Integrar empuje constante durante el intervalo.

\[\fitmath{\mathrm{IA}=\mathrm{FA}\cdot 120}\]

\[\fitmath{\mathrm{IA}=(0.4)\cdot 120\approx 48\,\mathrm{N\,s}}\]

\item Repetir el balance de momento y energía con la segunda velocidad.

\[\fitmath{\mathrm{FB}=\frac{2\cdot 600}{18000}\approx 0.06666667\,\mathrm{N}}\]

\item Calcular el segundo caudal.

\[\fitmath{\mathrm{mB}=\frac{2\cdot 600}{{18000}^{2}}\approx 3.703704\times10^{-6}\,\mathrm{kg/s}}\]

\item Impulso del segundo caso.

\[\fitmath{\mathrm{IB}=\mathrm{FB}\cdot 120}\]

\[\fitmath{\mathrm{IB}=(0.06666666667)\cdot 120\approx 8\,\mathrm{N\,s}}\]

\end{enumerate}

\end{studiumworked}
\figura{c01}{A potencia de chorro fija, aumentar la rapidez reduce el empuje. Modelo axial monocinético, sin término de presión.}
\section{Ejercicios propuestos}
\begin{enumerate}
\item Explique por qué un cohete funciona en vacío y por qué no puede prescindir de masa de reacción en el modelo descrito.
\item Un empuje constante de $0.12\,\mathrm N$ actúa durante $250\,\mathrm s$. Halle el impulso y distíngalo de un incremento de velocidad.
\item A potencia de chorro constante, duplique $v_e$. Prediga las razones de empuje, caudal e impulso específico antes de calcular.
\item Compare dos presupuestos de masa: uno que incluye sólo propulsante y otro que incluye fuente eléctrica. ¿Qué conclusión puede cambiar?
\end{enumerate}
\resumen{El sistema intercambia momento con materia expulsada. Empuje, impulso total e impulso específico no son intercambiables. Una elevada velocidad efectiva ahorra propulsante en ciertos modelos, pero exige energía y puede reducir el empuje. La elección depende de la misión y de una frontera de sistema explícita.}


\studiumchaptercheck
